% ------------------------------------------------------------------
%  Congedo — Sulle spalle dei giganti  (capitolo non numerato)
%  Ricerca di riferimento: ricerca 03 §4.2 (Ugo), §4.3 (Bernardo di
%    Chartres); ricerca 01 §17.1 (dodici principi); cap. 1 (Seneca)
%  Scheda: ricerca/06_indice-ragionato.md, §6.2
%  Epigrafe: ✔ verificata sul testo (Didascalicon VI, 3, The Latin Library)
%
%  STATO: prima stesura (07/10/2026), circa 2.200 parole (la scheda ne
%  prevedeva 1.500: lo scarto è quasi tutto nei dodici principi e nella
%  prova finale, che sono elenchi). Fonti verificate il 07/10/2026
%  (registro 00, K-28, K-37, K-38, S-121, S-122; note in
%  ricerca/fonti-scaricate-2026-10-07/_verifiche-online-2026-10-07.md):
%  - Ugo, Didascalicon VI, 3: ricordi di scuola (nomi delle cose,
%    sassolini, triangoli col carbone, corda tesa sul legno), «haec
%    puerilia quidem fuerant, sed tamen non inutilia, neque ea nunc
%    scire stomachum meum onerat»; «omnia disce… coartata scientia
%    iucunda non est» (The Latin Library).
%  - Giovanni di Salisbury, Metalogicon I, 24 («cogebantur exsolvere
%    singuli die sequenti aliquid eorum quae praecedenti audierant…
%    erat enim apud eos praecedentis discipulus sequens dies») e III, 4
%    (nani e giganti), testo della Patrologia Latina 199 (Logic Museum);
%    la PL stampa «incidentes», le edizioni critiche «insidentes».
%  - Bernardo: alla scuola di Chartres dal 1115 circa, cancelliere fino
%    al 1124; Giovanni lo conobbe attraverso i suoi allievi (◐).
%  - Vetrate del transetto sud di Chartres, circa 1221: evangelisti
%    sulle spalle dei profeti (Friends of Chartres).
%  - Newton a Hooke, 5 febbraio 1675 secondo il calendario inglese
%    (1676 secondo il nostro): Historical Society of Pennsylvania,
%    Linda Hall Library.
%  Nessun dato numerico nuovo; nessun aneddoto personale inventato
%  (la dedica alla maestra di italiano è già nel frontespizio).
% ------------------------------------------------------------------
\capitolosenzanumero{Sulle spalle dei giganti}
\label{cap:congedo}

\epigrafe[Impara tutto: vedrai poi che nulla è superfluo.]{Omnia disce, videbis postea nihil esse superfluum.}{Ugo di San Vittore, Didascalicon VI, 3}

\iniziale{V}{erso la fine} del \emph{Didascalicon}, il manuale dello
studente che hai incontrato nel capitolo~\ref{cap:storia}, Ugo di San
Vittore smette per un momento di dare regole e racconta di sé. Quando
era ancora uno scolaro, scrive, si era messo d'impegno a imparare il
nome di tutte le cose che gli capitavano sotto gli occhi o tra le
mani, perché non può indagare la natura delle cose chi non ne conosce
ancora i nomi. Disponeva sassolini per fare i conti; disegnava sul
pavimento, con pezzi di carbone, i triangoli, per vedere con i propri
occhi in che cosa differiscono quello ottuso, quello retto e quello
acuto; faceva scorrere un ponticello sotto una corda tesa su un legno,
per sentire con l'orecchio la differenza tra i suoni. \enquote{Erano
cose da ragazzi,} conclude, \enquote{ma non inutili; e saperle, adesso,
non mi pesa sullo stomaco}. Poi il consiglio che apre queste pagine:
impara tutto, vedrai poi che nulla è superfluo. \emph{Coartata scientia
iucunda non est}: un sapere ristretto non dà gioia.\footnote{Ugo di San
Vittore, \emph{Didascalicon} VI, 3 (circa 1128). L'immagine dello
stomaco non è casuale: per i lettori medievali un testo si
\emph{ruminava} (capitolo~\ref{cap:storia}), e ciò che è stato
digerito bene non pesa.}

Questo libro finisce dove finiva il suo: con l'invito a imparare molto,
e bene, e con la promessa che ciò che si impara bene non va perduto.
Prima di salutarti, però, vorrei raccontarti di un altro maestro dello
stesso secolo, e di una frase che forse conosci già.

\section*{Nani e giganti}

Negli anni in cui Ugo insegnava a Parigi, la scuola della cattedrale
di Chartres era guidata da un maestro di grammatica, Bernardo, che ne
fu cancelliere fino al 1124. Di suo non ci è rimasto quasi nulla. Lo
conosciamo soprattutto attraverso Giovanni di Salisbury, un inglese
venuto a studiare in Francia qualche anno più tardi, quando Bernardo
non insegnava più, e formato da alcuni dei suoi allievi. Nel 1159
Giovanni scrisse un libro in difesa degli studi, il \emph{Metalogicon},
e vi descrisse la scuola di Bernardo come un modello.

Il ritratto ha un dettaglio che, arrivato a questo punto del libro,
riconoscerai subito. Ogni giorno, racconta Giovanni, ciascun allievo
era tenuto a restituire qualcosa di ciò che aveva ascoltato il giorno
prima, chi di più, chi di meno: \enquote{presso di loro, infatti, il
giorno seguente era allievo del precedente}.\footnote{Giovanni di
Salisbury, \emph{Metalogicon} I, 24: \emph{cogebantur exsolvere singuli
die sequenti aliquid eorum quae praecedenti audierant, alii plus, alii
minus: erat enim apud eos praecedentis discipulus sequens dies}. Il
metodo di Bernardo aveva anche un lato che non rimpiangiamo: chi non
si applicava era spinto, secondo i casi, con ammonimenti oppure con la
frusta e con le punizioni.} È la sentenza di Publilio Siro che apre il
capitolo~\ref{cap:transizione}; ed è, otto secoli prima che qualcuno la
misurasse in laboratorio, la pratica del recupero
(capitolo~\ref{cap:recupero}): richiamare ogni giorno ciò che si è
imparato il giorno prima.

Nello stesso libro Giovanni riporta il detto per cui Bernardo è
ricordato ancora oggi. \enquote{Bernardo di Chartres diceva che noi
siamo come nani seduti sulle spalle di giganti, così da poter vedere
più cose di loro e più lontane, non per l'acutezza della nostra vista o
per l'altezza del nostro corpo, ma perché siamo sollevati e portati in
alto dalla grandezza dei giganti.}\footnote{Giovanni di Salisbury,
\emph{Metalogicon} III, 4: \emph{Dicebat Bernardus Carnotensis nos esse
quasi nanos gigantium humeris insidentes, ut possimus plura eis et
remotiora videre}… (alcune edizioni stampano \emph{incidentes}). I
giganti, per Bernardo, erano gli autori antichi che si leggevano nelle
scuole.}

L'immagine ha avuto una lunga fortuna. Nelle grandi vetrate del
transetto sud della cattedrale di Chartres, realizzate intorno al 1221,
i quattro evangelisti siedono sulle spalle di quattro profeti molto più
grandi di loro. E nel 1676 Isaac Newton, in una lettera a un collega
con cui era in polemica, scrisse che se aveva visto più lontano era
perché stava sulle spalle di giganti: una frase che spesso gli viene
attribuita come sua, e che aveva già più di cinque secoli.\footnote{La
lettera, a Robert Hooke, è datata 5 febbraio 1675 secondo il calendario
inglese del tempo, in cui l'anno cominciava a marzo: per noi è il
1676.}

Il detto di Bernardo ha due metà, ed entrambe contano. La prima è una
lezione di modestia: se vediamo più lontano degli antichi, non è perché
siamo più intelligenti di loro. La seconda è una promessa: vediamo
davvero più lontano. Ma c'è una condizione, che la frase dà per
scontata. I giganti non sollevano chi resta seduto ai loro piedi;
sollevano chi ci sale sopra. E salire, per Bernardo e per Ugo, voleva
dire una cosa precisa: imparare, ricordare, portare in sé ciò che i
giganti avevano scritto.

Oggi i giganti sono tutti a portata di mano. Qualunque libro è a un
clic di distanza, e un'intelligenza artificiale te ne riassume uno in
pochi secondi (capitolo~\ref{cap:ia}). Ma un gigante che hai soltanto
consultato non ti porta da nessuna parte: si vede lontano solo dalle
spalle su cui si è saliti. Ciò che sai diventa gli occhi con cui
guardi: è il senso più profondo del capitolo~\ref{cap:capire}, dove
abbiamo visto che si pensa con ciò che si sa. \emph{Omnia disce} non
vuol dire imparare a memoria tutto e qualunque cosa; vuol dire non
disprezzare nulla di ciò che regge il resto, nemmeno i nomi delle cose
e i triangoli disegnati sul pavimento.

\section*{I giganti di questo libro}

Ne hai incontrati molti, in queste pagine. Simonide, che ricordava il
posto di ciascun ospite intorno alla tavola; Aristotele, per cui tutti
gli uomini desiderano sapere; Quintiliano, Agostino, Ugo, Tommaso;
Comenio, che voleva insegnare tutto a tutti; Ebbinghaus, chino per anni
sulle sue sillabe senza senso; William James. E poi, più vicini a noi,
gli studiosi che hanno messo alla prova ciò che gli antichi avevano
intuito: Bahrick con il suo spagnolo dimenticato e ritrovato, Robert e
Elizabeth Bjork con le loro difficoltà desiderabili, Roediger e
Karpicke con i loro test, e molti altri i cui nomi riempiono la
bibliografia. Accanto a loro ci sono giganti senza nome: le migliaia di
studenti, quasi sempre universitari, che si sono prestati agli
esperimenti da cui viene gran parte di ciò che sai ora sullo studio.

Poi ci sono i giganti che nessuna bibliografia cita: chi ti ha
insegnato a leggere, a fare di conto, a scrivere un pensiero in ordine.
Ognuno ha i suoi. Questo libro è dedicato a una di loro.

E c'è un'ultima cosa, che il detto di Bernardo lascia in ombra. I
giganti non sono nati giganti: lo sono diventati studiando, e qualcuno
è salito sulle loro spalle mentre ancora salivano. Ciò che impari bene
non resta soltanto tuo. Lo spiegherai a un compagno, e nello spiegarlo
lo imparerai meglio (capitolo~\ref{cap:insegnare}); forse un giorno lo
insegnerai, in un'aula o a un figlio; di certo lo userai per qualcuno.
Un giorno, per qualcuno, sarai tu il gigante.

\section*{Ritorno a Seneca}

Il libro è cominciato con una frase rovesciata. Seneca scriveva
\emph{non vitae sed scholae discimus}, impariamo per la scuola e non
per la vita, e non era un augurio ma un lamento: si studia per la
prova, per la figura, per la gara, e poi si dimentica. Seneca non
disprezzava lo studio; disprezzava lo studio che si esaurisce nella
prova, e chiedeva un sapere che rendesse migliori, non soltanto più
dotti (capitolo~\ref{cap:vita}).

In ventisei capitoli abbiamo provato a rovesciare la sua frase una
seconda volta, con i fatti. Studiare per la vita non è un motto da
scrivere sopra la porta di una scuola: è il risultato di un metodo, e
il metodo è noto. Richiamare invece di rileggere; distanziare invece di
concentrare; alternare; spiegare; collegare; e sapere perché. Ciò che
si impara così non svanisce dopo l'appello: entra in quella parte della
memoria che Bahrick chiamava \emph{permastore} e diventa tuo, come il
cibo digerito di cui parlava Seneca stesso in un'altra lettera
(capitolo~\ref{cap:capire}).

Nessun metodo di studio può rendere buoni: su questo Seneca aveva
ragione, e il libro non lo promette. Ma ciò che sai bene cambia ciò che
vedi, e ciò che puoi fare per gli altri. È per questo che, alla domanda
\enquote{a che cosa mi serve questo?}, il libro ti ha chiesto ogni
volta di aggiungere: \enquote{e per chi?}.

\section*{Il metodo in dodici principi}

La scheda~\ref{sch:metodo} riassume \emph{che cosa}
fare. Qui, per l'ultima volta, c'è il \emph{perché}: i dodici principi
su cui il libro si regge, ciascuno con i capitoli in cui è spiegato e
messo alla prova.

\begin{enumerate}[leftmargin=1.8em, itemsep=0.35\baselineskip]
  \item \textbf{Imparare è cambiare la memoria a lungo termine.} Se lì
        nulla è cambiato, non hai imparato, anche se hai capito.
        {\footnotesize\itshape (capp.~\ref{cap:cervello}--\ref{cap:dueforze})}
  \item \textbf{Ricordare fa ricordare.} Richiamare dalla memoria è il
        modo più potente di studiare; una prova serve a imparare, non
        solo a giudicare.
        {\footnotesize\itshape (cap.~\ref{cap:recupero})}
  \item \textbf{Il tempo è un alleato.} Lo studio distribuito batte lo
        studio concentrato; ciò che vuoi ricordare a lungo va ripreso a
        intervalli crescenti.
        {\footnotesize\itshape (cap.~\ref{cap:tempo})}
  \item \textbf{La fatica giusta è un buon segno.} Le difficoltà
        desiderabili rallentano la prestazione di oggi e costruiscono
        l'apprendimento di domani, purché siano superabili.
        {\footnotesize\itshape (cap.~\ref{cap:dueforze})}
  \item \textbf{Non fidarti della sensazione di sapere.} Familiarità e
        scorrevolezza ingannano; solo un recupero onesto, a distanza di
        tempo, ti dice che cosa sai.
        {\footnotesize\itshape (cap.~\ref{cap:inganni})}
  \item \textbf{Mescola.} Alternare tipi di problemi ed esempi insegna
        a distinguerli e a scegliere la strada giusta.
        {\footnotesize\itshape (cap.~\ref{cap:alternanza})}
  \item \textbf{Chiediti perché.} Spiega a te stesso, collega al noto,
        prova a rispondere prima di ricevere la risposta; sbaglia senza
        paura e correggi subito.
        {\footnotesize\itshape (capp.~\ref{cap:capire} e~\ref{cap:insegnare})}
  \item \textbf{Si pensa con ciò che si sa.} La conoscenza è la base
        della comprensione e del pensiero critico; memorizzare e capire
        si sostengono a vicenda.
        {\footnotesize\itshape (capp.~\ref{cap:capire}--\ref{cap:immagini})}
  \item \textbf{Una cosa alla volta, ben guidata.} La memoria di lavoro
        è piccola: a chi comincia servono passi brevi ed esempi svolti,
        a chi è esperto più autonomia.
        {\footnotesize\itshape (capp.~\ref{cap:cervello} e~\ref{cap:docenti})}
  \item \textbf{Sapere perché studi cambia come studi.} Uno scopo
        personale, e il valore che vedi in ciò che studi, sostengono la
        fatica che imparare richiede.
        {\footnotesize\itshape (capp.~\ref{cap:sapere} e~\ref{cap:motivazione})}
  \item \textbf{Il corpo impara con la mente.} Il sonno prima di tutto;
        poi il movimento, le pause vere, l'attenzione protetta dalle
        distrazioni.
        {\footnotesize\itshape (capp.~\ref{cap:corpo}--\ref{cap:attenzione})}
  \item \textbf{Gli strumenti non pensano per te.} Libri, appunti,
        Internet e intelligenza artificiale aiutano quando ti
        costringono a pensare, e nuocciono quando pensano al posto tuo.
        {\footnotesize\itshape (cap.~\ref{cap:ia})}
\end{enumerate}

\section*{L'ultima prova}

Nel primo capitolo ti avevo chiesto di prendere un foglio e di
scrivere, senza aprire nulla, tutto ciò che ricordavi di un esame
superato un anno prima. Ora ti chiedo di fare lo stesso con questo
libro. Chiudilo, prendi un foglio e rispondi per iscritto alle domande
che seguono; soltanto dopo, controlla nei capitoli indicati in fondo.

Non preoccuparti di ciò che non ricordi: è la curva dell'oblio, e vale
anche per questo libro. Preoccupati piuttosto di cercare la risposta
prima di arrenderti, perché è in quello sforzo che il ricordo si
rafforza. Le domande sono mescolate di proposito. E se vuoi fare
l'ultima cosa che il libro ti chiede, ripeti la prova fra un mese.

\begin{riquadro}{L'ultima prova}
\begin{enumerate}
  \item Perché la rilettura dà la sensazione di sapere, e come puoi
        accorgerti che quella sensazione ti sta ingannando?
  \item Hai un esame fra due mesi. Come distribuiresti le riprese di
        un argomento, e perché proprio così?
  \item Che differenza c'è tra prestazione e apprendimento? Fai un
        esempio di difficoltà desiderabile.
  \item Che cosa ha capito Simonide dopo il crollo della sala, e quando
        una mnemotecnica serve davvero?
  \item Perché, mescolando esercizi di tipo diverso, sembra di imparare
        di meno mentre si impara di più?
  \item Perché memorizzare e capire non sono in contrasto?
  \item Che cosa fanno alla memoria una notte di sonno e un telefono
        sul tavolo?
  \item Quando l'intelligenza artificiale ti aiuta a imparare, e quando
        pensa al posto tuo?
  \item Perché sapere a che cosa serve ciò che studi cambia il modo in
        cui lo studi?
  \item Che cosa dice la ricerca sugli \enquote{stili di
        apprendimento}?
  \item Per chi insegna: come si può fare di un esame un'occasione per
        imparare, e non soltanto per misurare?
  \item Che cosa diceva Seneca, davvero, sullo studio \enquote{per la
        scuola}?
  \item Chi diceva che siamo nani sulle spalle di giganti, e che cosa
        avevano in comune la sua scuola e la pratica del recupero?
\end{enumerate}

\smallskip
{\footnotesize\itshape Dove trovare le risposte: (1)
cap.~\ref{cap:inganni}; (2) cap.~\ref{cap:tempo} e
scheda~\ref{sch:calendario}; (3) cap.~\ref{cap:dueforze}; (4)
capp.~\ref{cap:storia} e~\ref{cap:immagini}; (5)
cap.~\ref{cap:alternanza}; (6) cap.~\ref{cap:capire}; (7)
capp.~\ref{cap:corpo} e~\ref{cap:attenzione}; (8) cap.~\ref{cap:ia};
(9) capp.~\ref{cap:sapere} e~\ref{cap:motivazione}; (10)
cap.~\ref{cap:miti}; (11) cap.~\ref{cap:valutare}; (12)
cap.~\ref{cap:vita}; (13) questo congedo.\par}
\end{riquadro}

\asterismo

Il primo capitolo finiva con un invito: si comincia chiudendo questo
libro e provando a ricordare. Il libro finisce allo stesso modo, perché
il metodo è lo stesso dall'inizio alla fine, e perché nessun libro, da
solo, ti farà salire sulle spalle di nessuno. Puoi farlo solo tu, un
gradino alla volta, ogni giorno allievo del giorno prima.

I giganti sono lì, e le loro spalle sono larghe. Buono studio; e,
proprio per questo, buona vita.

\finecapitolo
